\chapter{Projection-Free Nonsmooth Optimization with Functional Constraints}
\label{Section:projfree}

\noindent\textbf{Source and author contributions.}
This chapter is based on the algorithm published jointly with
Michael J. Neely~\cite{asgari2024nonsmooth}.
Neely suggested extending the original projection-free method to functional
constraints and directed me to the work of Neely and
Yu~\cite{neely2019lagrangian}. Their treatment of Lagrangian methods served
as a starting point and source of inspiration for my development of the
multiplier update in the published paper. I proved the guarantees of the
resulting method. Neely reviewed the method and analysis.
The original update rules remain unchanged; only the constraints are replaced by their positive parts. GPT-6 Astra suggested this modification, developed the corresponding proofs, and helped with the chapter’s exposition.

\section{Problem and method}

This chapter extends the general convex method in
Section~\ref{sec:pfdom:convex} to minimization of $f$ over $\X$
subject to $h_i(x)\le0$, $i=1,\ldots,m$.
Write $h:\Real^n\to\Real^m$ for the constraint map,
$h(x)\coloneqq(h_1(x),\ldots,h_m(x))^\top$.
Here $\X$ is compact and convex, and the objective and constraints may be
nonsmooth. Linear optimization over $\X$ is assumed inexpensive, whereas
Euclidean projection may be costly.

The linear oracle exploits the structure of $\X$. Linear minimization
over the intersection $\{x\in\X:h(x)\le0\}$ may be more costly, and no
inexpensive oracle over this intersection is assumed. The algorithm therefore
retains linear minimization over $\X$ and handles the functional
inequalities through multiplier updates, allowing controlled violation.
Its returned point belongs to $\X$ exactly, with bounds on its objective
gap and functional-constraint violation.

\section{Related work}
\label{sec:ch4-related-work}

\subsection{Primal--dual and penalty methods}

The constraint extrapolation method, ConEx, uses linearized constraints
\cite[Section 2.1]{BoobDengLan2023ConEx}. It covers convex nonsmooth
objectives and constraints and permits noisy constraint values and
subgradients. Its multiplier update avoids projection onto a bounded set.
Its primal step still requires a proximal operation over the original domain.

Exact penalty methods incorporate constraint violation into the objective
\cite{di1994exact}. Under the assumption that a Lagrange multiplier exists, a
sufficiently large coefficient makes the penalized and constrained problems
have the same minimizers. Choosing a coefficient that guarantees this
property a priori requires a suitable bound on the norm of a Lagrange
multiplier, and a larger coefficient can increase the Lipschitz bound used
in the analysis.

\subsection{Projection-free methods with functional constraints}

Lan, Romeijn, and Zhou combine constraint extrapolation with conditional
gradients in CoexCG and CoexDurCG \cite{lan2021conditional}. Their nonsmooth
extension constructs smooth approximations from explicit max
representations of the functions, requiring more information than
subgradients alone.

The Level Conditional Gradient method, LCG, solves a sequence of
subproblems that seek a point satisfying a prescribed set of inequalities,
using conditional gradients
\cite{cheng2026functional}. Its nonsmooth extension also requires max
representations and smoothing.

\subsection{Oracle requirements and iteration bounds}

The algorithm of Asgari and Neely~\cite{asgari2024nonsmooth} treats general
nonsmooth convex objectives and functional constraints. The published
paper permits stochastic subgradients and inexact linear minimization,
while constraint values are exact. It uses linear optimization over
$\X$ and projection onto an auxiliary set $\Y\supseteq\X$.
The exact deterministic specialization used here is compared with the
published updates in Section~\ref{sec:ch4-published-updates}.

For general convex nonsmooth problems,
ConEx~\cite[Theorem~2.3]{BoobDengLan2023ConEx} and the method of
Asgari and Neely~\cite[Theorems~1--2]{asgari2024nonsmooth} give
$O(\epsilon^{-2})$ iteration bounds, with constraint violation bounds depending
on the norm of a Lagrange multiplier. For structured nonsmooth problems,
CoexCG and CoexDurCG have the same $O(\epsilon^{-2})$ iteration order,
and their constraint violation bounds also depend on this norm
\cite[Corollaries~2.7 and~3.7]{lan2021conditional}.
LCG gives $O(\epsilon^{-2}\log(1/\epsilon))$ total inner iterations
for structured nonsmooth problems, with a bound independent of the
norm of a Lagrange multiplier
\cite[Theorem~2.4]{cheng2026functional}.
These bounds concern both objective accuracy and constraint violation, with
guarantees interpreted in expectation for stochastic variants.
Each iteration order above holds under the problem and oracle hypotheses
of the theorem or corollary cited with that result, with all constants
other than $\epsilon$ held fixed.
The cost of an iteration differs between methods, so these bounds alone
do not compare running time.

\section{Contributions and scope}

This chapter gives a combined proof of a joint bound on the objective gap
and squared constraint violation at the averaged output for the algorithm
of Asgari and Neely~\cite{asgari2024nonsmooth}, applied to clipped
constraints with exact deterministic oracles.
The method uses subgradients and constraint values without requiring
smoothness, a special smoothing
representation, or proximal operators for the objective and constraint
functions. All update coefficients depend on the query budget $T$
and the known bounds $D,L,G$.

The objective gap has an $O(T^{-1/2})$ upper bound. Under
Assumption~\ref{assum:Lagrange multiplier}, the constraint violation norm
also has an $O(T^{-1/2})$ upper bound, as stated in
\eqref{eq:ch4-cb-violation-theorem}.

The method uses at most $T$ calls to each oracle: $T$ constraint-vector
evaluations, $T-1$ joint subgradient queries, and $T-1$ linear-oracle calls.
It also performs $T-1$ projections onto the auxiliary set.
These counts assume that $x_1\in\X$ is supplied; obtaining $x_1$ with one
additional linear-oracle call keeps the linear-oracle count within the
budget $T$.
Exact constraint values $h(y)$ and subgradient selections satisfying
\eqref{eq:ch4-cb-oracles} must be available at every $y\in\Y$.
The published paper's stochastic and inexact-oracle results,
application examples, and numerical experiments are outside the scope of
the present chapter.

\section{Problem formulation and oracle assumptions}
\label{sec:ch4-assumptions}

We follow the notation of Section~\ref{sec:intro:notation}. All norms are
Euclidean; vector inequalities and positive parts are componentwise. Let
$f,h_1,\ldots,h_m:\Real^n\to\Real$ be finite-valued and convex. We consider

\begin{equation}\label{eq:ch4-cb-problem}
f^*\coloneqq\min_{x\in\X}\{f(x):h(x)\le0\},
\qquad h(x)\coloneqq(h_1(x),\ldots,h_m(x))^\top.
\end{equation}

\begin{Thm:Assumption}[Domain and exact deterministic oracles]\label{ass:ch2-oracles}
Assume $\X$ is nonempty, compact, and convex; $\Y$ is closed and convex; and a feasible minimizer $x^*$ exists. Fix positive bounds $D,L,G$ such that

\begin{equation}\label{eq:ch4-cb-domain}
\begin{gathered}
\X\subseteq\Y,\qquad \sup_{x,x'\in\X}\|x-x'\|\le D,\\
x^*\in\X,\qquad h(x^*)\le0,\qquad f(x^*)=f^*.
\end{gathered}
\end{equation}

The constants $D,L,G$ are known positive upper bounds.
For every $v\in\Real^n$, an exact linear oracle returns a minimizer of
$\langle v,x\rangle$ over $x\in\X$, and the Euclidean projection
$\Pi_\Y(v)$ is available.
At every $y\in\Y$, the exact oracle supplies the constraint values $h(y)$
and subgradient selections satisfying
\begin{equation}\label{eq:ch4-cb-oracles}
\begin{aligned}
s(y)&\in\partial f(y),&\quad \|s(y)\|&\le L,\\
g_i(y)&\in\partial h_i(y),&\quad \|g_i(y)\|&\le G_i,
\qquad \sum_{i=1}^mG_i^2\le G^2.
\end{aligned}
\end{equation}

The constants $G_i\ge0$ need not be individually known.
\end{Thm:Assumption}

For computational efficiency, choose $\Y$ with an inexpensive projection,
such as an enclosing ball or box. The bounds $L$ and $G$ must hold
throughout that set.
Section~\ref{sec:pfdom:auxiliary-choice} explains how enlarging the query
set can increase these bounds.

Define the clipped constraints and their oracle selections by
\begin{equation}\label{eq:ch4-cb-clipped-constraints}
c_i(y)\coloneqq[h_i(y)]_+=\max\{h_i(y),0\},\qquad
c(y)\coloneqq(c_1(y),\ldots,c_m(y))^\top,\qquad y\in\Real^n.
\end{equation}
\begin{equation}\label{eq:ch4-cb-clipped-subgradients}
k_i(y)\coloneqq
\begin{cases}
g_i(y),&h_i(y)>0,\\
0,&h_i(y)\le0,
\end{cases}
\qquad y\in\Y.
\end{equation}

\begin{Thm:Lemma}[Clipped constraints and Lipschitz bounds]\label{lem:ch4-cb-lipschitz}
Under Assumption~\ref{ass:ch2-oracles}, each $c_i$ is convex and
\begin{equation}\label{eq:ch4-cb-clipped-zero-violation-set}
\{x\in\X:h(x)\le0\}
=\{x\in\X:c(x)=0\}
=\{x\in\X:c(x)\le0\}.
\end{equation}
At every $y\in\Y$,
\begin{equation}\label{eq:ch4-cb-clipped-oracles}
k_i(y)\in\partial c_i(y),\qquad
\|k_i(y)\|\le G_i,\qquad
\sum_{i=1}^m\|k_i(y)\|^2\le G^2.
\end{equation}
On $\Y$, $f$ is $L$-Lipschitz, and $h$, $c$, and $y\mapsto\|c(y)\|$
are $G$-Lipschitz. The function $y\mapsto\|c(y)\|^2$ is convex on
$\Real^n$.
\end{Thm:Lemma}
\begin{proof}
Each $c_i=\max\{h_i,0\}$ is convex as the maximum of two convex
functions. Since $c_i\ge0$ and $c_i(x)=0$ exactly when $h_i(x)\le0$,
the feasible sets in \eqref{eq:ch4-cb-clipped-zero-violation-set} coincide.

Fix $y\in\Y$ and $x\in\Real^n$. If $h_i(y)>0$, then
\begin{equation}\label{eq:ch4-cb-clipped-positive-subgradient}
\begin{aligned}
c_i(x)&\ge h_i(x)
\ge h_i(y)+\langle g_i(y),x-y\rangle\\
&=c_i(y)+\langle k_i(y),x-y\rangle.
\end{aligned}
\end{equation}
The inequalities use $c_i\ge h_i$ and the subgradient inequality for
$h_i$. If $h_i(y)\le0$, including $h_i(y)=0$, then $c_i(y)=0$ is a
global minimum of $c_i$, so the choice $k_i(y)=0$ is a subgradient.
Thus $k_i(y)\in\partial c_i(y)$ in both cases. Its norm bounds in
\eqref{eq:ch4-cb-clipped-oracles} follow from
$k_i(y)\in\{g_i(y),0\}$ and \eqref{eq:ch4-cb-oracles}.

For $x,y\in\Y$,
\begin{equation}\label{eq:ch4-cb-scalar-lipschitz}
\begin{aligned}
|f(x)-f(y)|&\le L\|x-y\|,\\
|h_i(x)-h_i(y)|&\le G_i\|x-y\|,\\
|c_i(x)-c_i(y)|&\le G_i\|x-y\|.
\end{aligned}
\end{equation}
These inequalities use Lemma~\ref{lem:intro:lipschitz} with
\eqref{eq:ch4-cb-oracles} and \eqref{eq:ch4-cb-clipped-oracles}.
Squaring and summing the coordinate bounds, using
$\sum_iG_i^2\le G^2$, and taking square roots gives
\begin{equation}\label{eq:ch4-cb-lipschitz}
\begin{aligned}
\|h(x)-h(y)\|&\le G\|x-y\|,\\
\|c(x)-c(y)\|&\le G\|x-y\|.
\end{aligned}
\end{equation}
The reverse triangle inequality then gives
\begin{equation}\label{eq:ch4-cb-v-lipschitz}
\begin{aligned}
\bigl|\|c(x)\|-\|c(y)\|\bigr|
&\le\|c(x)-c(y)\|\le G\|x-y\|.
\end{aligned}
\end{equation}
Finally, for $x,y\in\Real^n$ and $0\le\lambda\le1$,
\begin{equation}\label{eq:ch4-cb-squared-violation-convexity}
\begin{aligned}
\|c(\lambda x+(1-\lambda)y)\|^2
&\overset{\mathrm{(A)}}{\le}
\|\lambda c(x)+(1-\lambda)c(y)\|^2\\
&\overset{\mathrm{(B)}}{\le}
\lambda\|c(x)\|^2+(1-\lambda)\|c(y)\|^2.
\end{aligned}
\end{equation}
Here (A) uses convexity and nonnegativity of each $c_i$ and monotonicity
of the squared norm on the nonnegative orthant; (B) uses convexity of the
squared norm.
\end{proof}

\section{Projection-free algorithm}
\label{sec:ch2-development}

The method maintains a domain point $x_t\in\X$ and a query point
$y_t\in\Y$. Linear optimization produces the $x_t$, while
subgradients and constraint values are evaluated at the $y_t$.
The equality multiplier $q_t\in\Real^n$ couples the two sequences.
The inequality multiplier $w_t\in\Real_+^m$ is updated using the
linearized clipped constraints.

Let $T\ge1$ be an integer query budget for each oracle. We count a full
constraint-vector evaluation as one query and a joint evaluation of the
objective and constraint subgradients as one subgradient query. Oracle
counts are reported separately by type, including any linear-oracle call
used to obtain the initial point.
For problem \eqref{eq:ch4-cb-problem}, use the update coefficients
\begin{equation}\label{eq:ch4-cb-tuning}
a\coloneqq\frac{2L\sqrt T}{D},\qquad
\beta\coloneqq\frac{a}{2G^2},\qquad
\eta\coloneqq\frac aT.
\end{equation}
Initialize
\begin{equation}\label{eq:ch4-cb-init}
y_1\coloneqq x_1\in\X,\qquad q_1\coloneqq0,\qquad w_1\coloneqq0.
\end{equation}
Evaluate $h(y_1)$ and form $c(y_1)$. For $t=1,\ldots,T-1$, use
\begin{equation}\label{eq:ch4-cb-oracle-data}
s_t\coloneqq s(y_t),\qquad
K_t\coloneqq\begin{bmatrix}k_1(y_t)^\top\\ \vdots\\k_m(y_t)^\top\end{bmatrix}.
\end{equation}
Define the affine model of the clipped constraints by
\begin{equation}\label{eq:ch4-cb-affine-model}
\widehat c_t(y)\coloneqq c(y_t)+K_t(y-y_t).
\end{equation}
Then update
\begin{equation}\label{eq:ch4-cb-linear-oracle}
x_{t+1}\in\arg\min_{x\in\X}\langle q_t,x\rangle,
\end{equation}
\begin{equation}\label{eq:ch4-cb-primal-update}
y_{t+1}\coloneqq\Pi_{\Y}\left(
\frac{a y_t+\eta x_{t+1}+q_t-s_t-\beta K_t^\top(w_t+c(y_t))}{a+\eta}
\right),
\end{equation}
\begin{equation}\label{eq:ch4-cb-multiplier-updates}
\begin{aligned}
q_{t+1}&\coloneqq q_t-\eta(y_{t+1}-x_{t+1}),\\
w_{t+1}&\coloneqq[w_t+\widehat c_t(y_{t+1})]_+.
\end{aligned}
\end{equation}
Evaluate $h(y_{t+1})$ and form $c(y_{t+1})$ after the update.
Equations~\eqref{eq:ch4-cb-clipped-constraints},
\eqref{eq:ch4-cb-init}, and \eqref{eq:ch4-cb-multiplier-updates} give
\begin{equation}\label{eq:ch4-cb-nonnegative-weights}
w_t\ge0,\qquad c(y_t)\ge0,\qquad w_t+c(y_t)\ge0.
\end{equation}

The method returns $\bar x_T$, defined below. We denote its
objective gap and constraint violation by $\Delta_T$ and $V_T$:
\begin{equation}\label{eq:ch4-cb-output}
\begin{gathered}
\bar x_T\coloneqq\frac1T\sum_{t=1}^Tx_t,\\
\Delta_T\coloneqq f(\bar x_T)-f^*,\qquad
V_T\coloneqq\|c(\bar x_T)\|=\|[h(\bar x_T)]_+\|.
\end{gathered}
\end{equation}

Convexity of $\X$ ensures that $\bar x_T\in\X$. The functional
constraints need not hold exactly at this average; their violation is
measured by $V_T$. The objective gap $\Delta_T$ can therefore be negative
when the output is infeasible.

Algorithm~\ref{alg:functional} collects
\eqref{eq:ch4-cb-tuning}--\eqref{eq:ch4-cb-output}.
\begin{algorithm}[H]
\caption{Projection-free method with functional constraints}
\label{alg:functional}
\begin{algorithmic}[1]
\Require Integer query budget $T\ge1$, bounds $D,L,G>0$, and $x_1\in\X$
\State Set $a,\beta,\eta$ by \eqref{eq:ch4-cb-tuning}
\State $y_1\gets x_1$, $q_1\gets0$, $w_1\gets0$
\State Evaluate $h(y_1)$ and form $c(y_1)$
\For{$t=1,\ldots,T-1$}
    \State Query $(s_t,g_1(y_t),\ldots,g_m(y_t))$ jointly
    \State Form $K_t$ using \eqref{eq:ch4-cb-clipped-subgradients} and \eqref{eq:ch4-cb-oracle-data}
    \State $x_{t+1}\gets\operatorname{LOO}_\X(q_t)$
    \State $\displaystyle y_{t+1}\gets\Pi_\Y\left(
    \frac{a y_t+\eta x_{t+1}+q_t-s_t-\beta K_t^\top(w_t+c(y_t))}{a+\eta}\right)$
    \State $q_{t+1}\gets q_t-\eta(y_{t+1}-x_{t+1})$
    \State $w_{t+1}\gets[w_t+\widehat c_t(y_{t+1})]_+$ using \eqref{eq:ch4-cb-affine-model}
    \State Evaluate $h(y_{t+1})$ and form $c(y_{t+1})$
\EndFor
\State \Return $\displaystyle\bar x_T=T^{-1}\sum_{t=1}^T x_t$
\end{algorithmic}
\end{algorithm}

\subsection{Relation to the published updates}
\label{sec:ch4-published-updates}

This chapter applies the published algorithm~\cite{asgari2024nonsmooth}
to $c_i=[h_i]_+$ with exact deterministic oracles.
Relative to the published notation, $q_t=-\eta Q_t$, $w_t=W_t$, and
$a=\alpha+2G^2\beta$. Since $c_i\ge0$, the published initialization
becomes $w_1=[-c(y_1)]_+=0$, and the lower bound
$[-c(y_{t+1})]_+$ in its multiplier update is also zero. This gives
\eqref{eq:ch4-cb-primal-update}--\eqref{eq:ch4-cb-multiplier-updates}.

The tuning in \eqref{eq:ch4-cb-tuning} sets $a=2G^2\beta>0$,
corresponding to $\alpha=0$ in the published notation. Although the
published analysis requires a separate positive $\alpha$, the proof in
Section~\ref{sec:ch4-cb-proof} bounds the objective-subgradient term by
$L\|y_{t+1}-y_t\|-(a/4)\|y_{t+1}-y_t\|^2\le L^2/a$;
see \eqref{eq:ch4-cb-square-bound} and \eqref{eq:ch4-cb-one-step}.
Thus the proof covers this parameter choice.

\section{Convergence analysis}
\label{sec:ch4-convergence}
\subsection{Main convergence guarantee}

\begin{Thm:Assumption}[Existence of a Lagrange multiplier]\label{assum:Lagrange multiplier}
There is a vector $\mu$ satisfying
\begin{equation}\label{eq:ch4-cb-multiplier}
\begin{gathered}
\mu\in\Real_+^m,\\
f(x)+\mu^\top h(x)\ge f^*\qquad(x\in\X).
\end{gathered}
\end{equation}
The multiplier $\mu$ is not an algorithm input.
\end{Thm:Assumption}
Assumption~\ref{assum:Lagrange multiplier} supplies the multiplier used in
the constraint violation bound and the objective lower bound in
\eqref{eq:ch4-cb-violation-theorem}. A sufficient condition for
Assumption~\ref{assum:Lagrange multiplier} is the existence of
$x^\circ\in\operatorname{relint}\X$ with $h_i(x^\circ)<0$ for every
$i=1,\ldots,m$ \cite{boyd_vandenberghe_2004}.

\begin{Thm:Theorem}[Objective gap and constraint violation bounds]\label{thm:ch4-cb-guarantee}
Under Assumption~\ref{ass:ch2-oracles}, Algorithm~\ref{alg:functional}
with \eqref{eq:ch4-cb-tuning} satisfies
\begin{equation}\label{eq:ch4-cb-theorem}
\frac{\Delta_T}{LD}
+\frac{\sqrt T}{4}\left(\frac{V_T}{GD}\right)^2
\le\frac3{\sqrt T},
\qquad \Delta_T\le\frac{3LD}{\sqrt T}.
\end{equation}
If Assumption~\ref{assum:Lagrange multiplier} also holds, then
\begin{equation}\label{eq:ch4-cb-violation-theorem}
V_T\le\frac{4GD}{\sqrt T}\left(1+\frac{G\|\mu\|}{L}\right),
\qquad -\|\mu\|V_T\le\Delta_T\le\frac{3LD}{\sqrt T}.
\end{equation}
\end{Thm:Theorem}

\subsection{Proof and oracle complexity}
\label{sec:ch4-cb-proof}

\begin{proof}[Proof of Theorem~\ref{thm:ch4-cb-guarantee}]
\postdisplaypenalty=10000
\interlinepenalty=10000

The primal and multiplier estimates give a telescoping bound on the
objective gap plus squared constraint violation at the query points.
The equality multiplier controls the distance between the query-point
average and the returned average, allowing this bound to be transferred
to the returned point in $\X$. The final step uses
Assumption~\ref{assum:Lagrange multiplier} to obtain a separate
constraint-violation bound.

\par\medskip\noindent\textbf{The inequality multiplier.}
For $t=1,\ldots,T-1$, write the query-point increment and the
parameter identities from \eqref{eq:ch4-cb-tuning} as
\begin{equation}\label{eq:ch4-cb-proof-notation}
\begin{gathered}
u_t\coloneqq y_{t+1}-y_t,\\
a=2\beta G^2=\eta T=\frac{2L\sqrt T}{D}.
\end{gathered}
\end{equation}
For $v\in\Real^n$, \eqref{eq:ch4-cb-oracle-data} gives
\begin{equation}\label{eq:ch4-cb-jacobian-bound}
\begin{aligned}
\|K_tv\|^2
&=\sum_i\langle k_i(y_t),v\rangle^2\\
&\overset{\mathrm{(A)}}{\le}\sum_i\|k_i(y_t)\|^2\|v\|^2
\overset{\mathrm{(B)}}{\le}G^2\|v\|^2,\\
\widehat c_t(x^*)&=c(y_t)+K_t(x^*-y_t)
\overset{\mathrm{(C)}}{\le}c(x^*)=0.
\end{aligned}
\end{equation}
\nopagebreak[3]
Here (A) is Cauchy--Schwarz, (B) uses
\eqref{eq:ch4-cb-clipped-oracles}, and (C) applies the subgradient inequality
to each $c_i$, using $k_i(y_t)\in\partial c_i(y_t)$ from
\eqref{eq:ch4-cb-clipped-oracles}.
The final equality uses \eqref{eq:ch4-cb-domain} and
\eqref{eq:ch4-cb-clipped-zero-violation-set}.

Expanding the inequality-multiplier update
\eqref{eq:ch4-cb-multiplier-updates} gives
\begin{equation}\label{eq:ch4-cb-W-drift}
\begin{aligned}
&\frac{\|w_{t+1}\|^2-\|w_t\|^2}{2}\\
&\quad=\frac12\|[w_t+\widehat c_t(y_{t+1})]_+\|^2-\frac12\|w_t\|^2\\
&\quad\overset{\mathrm{(A)}}{\le}
\frac12\|w_t+\widehat c_t(y_{t+1})\|^2-\frac12\|w_t\|^2\\
&\quad=w_t^\top\widehat c_t(y_{t+1})+\frac12\|c(y_t)+K_tu_t\|^2\\
&\quad=(w_t+c(y_t))^\top\widehat c_t(y_{t+1})
-\frac12\|c(y_t)\|^2+\frac12\|K_tu_t\|^2\\
&\quad\overset{\mathrm{(B)}}{\le}(w_t+c(y_t))^\top\widehat c_t(y_{t+1})
-\frac12\|c(y_t)\|^2+\frac{G^2}{2}\|u_t\|^2.
\end{aligned}
\end{equation}
\nopagebreak[3]
Here (A) uses $\|[v]_+\|\le\|v\|$; (B) uses
\eqref{eq:ch4-cb-jacobian-bound}.
The equalities use $\widehat c_t(y_{t+1})=c(y_t)+K_tu_t$.
Multiplying by $\beta>0$ and rearranging gives
\begin{equation}\label{eq:ch4-cb-constraint-work}
\begin{aligned}
&\beta(w_t+c(y_t))^\top\widehat c_t(y_{t+1})\\
&\quad\ge\frac\beta2\bigl(\|w_{t+1}\|^2-\|w_t\|^2\bigr)
+\frac\beta2\|c(y_t)\|^2-\frac{\beta G^2}{2}\|u_t\|^2.
\end{aligned}
\end{equation}

\par\medskip\noindent\textbf{The primal update.}
Completing the square in \eqref{eq:ch4-cb-primal-update} gives
\begin{equation}\label{eq:ch4-cb-primal-minimization}
\begin{aligned}
y_{t+1}=\arg\min_{y\in\Y}\biggl\{
&\langle s_t+\beta K_t^\top(w_t+c(y_t))-q_t,y\rangle\\
&+\frac\eta2\|y-x_{t+1}\|^2+\frac a2\|y-y_t\|^2\biggr\}.
\end{aligned}
\end{equation}
Apply \eqref{eq:pushback} in the Euclidean geometry
\eqref{eq:intro:quadratic-geometry}, with ambient space $\Real^n$,
feasible set $\Y$, reference point $y_t$,
coefficient $a$, and comparison point $x^*$. The convex function used in
\eqref{eq:pushback} is
\[
\psi_t(y)\coloneqq\langle s_t+\beta K_t^\top(w_t+c(y_t))-q_t,y\rangle
+\frac\eta2\|y-x_{t+1}\|^2.
\]
Expanding \eqref{eq:pushback} gives
\begin{equation}\label{eq:ch4-cb-quadratic-comparison}
\begin{aligned}
&\langle s_t+\beta K_t^\top(w_t+c(y_t))-q_t,y_{t+1}-x^*\rangle\\
&\quad+\frac\eta2\|y_{t+1}-x_{t+1}\|^2+\frac a2\|u_t\|^2\\
&\quad\overset{\mathrm{(A)}}{\le}
\frac\eta2\|x_{t+1}-x^*\|^2
+\frac a2\bigl(\|y_t-x^*\|^2-\|y_{t+1}-x^*\|^2\bigr)\\
&\quad\overset{\mathrm{(B)}}{\le}
\frac{\eta D^2}{2}
+\frac a2\bigl(\|y_t-x^*\|^2-\|y_{t+1}-x^*\|^2\bigr).
\end{aligned}
\end{equation}
\nopagebreak[3]
Here (A) is \eqref{eq:pushback} applied to
\eqref{eq:ch4-cb-primal-minimization}; (B) uses \eqref{eq:ch4-cb-domain}.

For the $q_t$ terms,
\begin{equation}\label{eq:ch4-cb-equality-identity}
\begin{aligned}
&-\langle q_t,y_{t+1}-x^*\rangle
+\frac\eta2\|y_{t+1}-x_{t+1}\|^2\\
&\quad=-\langle q_t,y_{t+1}-x_{t+1}\rangle
+\frac\eta2\|y_{t+1}-x_{t+1}\|^2
-\langle q_t,x_{t+1}-x^*\rangle\\
&\quad=\frac{\|q_t-\eta(y_{t+1}-x_{t+1})\|^2-\|q_t\|^2}{2\eta}
-\langle q_t,x_{t+1}-x^*\rangle\\
&\quad=\frac{\|q_{t+1}\|^2-\|q_t\|^2}{2\eta}
-\langle q_t,x_{t+1}-x^*\rangle\\
&\quad\overset{\mathrm{(A)}}{\ge}\frac{\|q_{t+1}\|^2-\|q_t\|^2}{2\eta}.
\end{aligned}
\end{equation}
\nopagebreak[3]
The equalities expand the square and use
\eqref{eq:ch4-cb-multiplier-updates}; (A) uses
\eqref{eq:ch4-cb-linear-oracle}.

Also,
\begin{equation}\label{eq:ch4-cb-subgradient-comparison}
\begin{aligned}
&\langle s_t,y_{t+1}-x^*\rangle\\
&\quad=\langle s_t,y_t-x^*\rangle+\langle s_t,u_t\rangle\\
&\quad\overset{\mathrm{(A)}}{\ge}f(y_t)-f^*+\langle s_t,u_t\rangle\\
&\quad\overset{\mathrm{(B)}}{\ge}f(y_t)-f^*-L\|u_t\|,\\
&(w_t+c(y_t))^\top K_t(y_{t+1}-x^*)\\
&\quad=(w_t+c(y_t))^\top\bigl(\widehat c_t(y_{t+1})-\widehat c_t(x^*)\bigr)\\
&\quad\overset{\mathrm{(C)}}{\ge}(w_t+c(y_t))^\top\widehat c_t(y_{t+1}).
\end{aligned}
\end{equation}
\nopagebreak[3]
The equalities use \eqref{eq:ch4-cb-affine-model} and
\eqref{eq:ch4-cb-proof-notation}.
Here (A) is the subgradient inequality with $f(x^*)=f^*$ from
\eqref{eq:ch4-cb-domain}; (B) is Cauchy--Schwarz with
\eqref{eq:ch4-cb-oracles}; (C) uses
\eqref{eq:ch4-cb-nonnegative-weights} and
\eqref{eq:ch4-cb-jacobian-bound}.

Substituting the multiplier identity \eqref{eq:ch4-cb-equality-identity}
and the subgradient bounds \eqref{eq:ch4-cb-subgradient-comparison}
into \eqref{eq:ch4-cb-quadratic-comparison} gives
\begin{equation}\label{eq:ch4-cb-primal-comparison}
\begin{aligned}
&f(y_t)-f^*+\beta(w_t+c(y_t))^\top\widehat c_t(y_{t+1})\\
&\quad+\frac a2\bigl(\|y_{t+1}-x^*\|^2-\|y_t-x^*\|^2\bigr)\\
&\quad+\frac{\|q_{t+1}\|^2-\|q_t\|^2}{2\eta}
+\frac a2\|u_t\|^2
\le\frac{\eta D^2}{2}+L\|u_t\|.
\end{aligned}
\end{equation}

Define
\begin{equation}\label{eq:ch4-cb-F-definition}
\begin{aligned}
F(y)&\coloneqq f(y)-f^*+\frac\beta2\|c(y)\|^2,\\
E_t&\coloneqq\frac a2\|y_t-x^*\|^2+\frac{\|q_t\|^2}{2\eta}+\frac\beta2\|w_t\|^2.
\end{aligned}
\end{equation}
The function $F$ is convex by \eqref{eq:ch4-cb-squared-violation-convexity}
and convexity of $f$. For $r\ge0$,
\begin{equation}\label{eq:ch4-cb-square-bound}
Lr-\frac a4r^2
=\frac{L^2}{a}-\frac a4\left(r-\frac{2L}{a}\right)^2
\overset{\mathrm{(A)}}{\le}\frac{L^2}{a}.
\end{equation}
\nopagebreak[3]
Here (A) drops the nonpositive term
$-(a/4)(r-2L/a)^2$.

Using \eqref{eq:ch4-cb-F-definition},
\begin{equation}\label{eq:ch4-cb-one-step}
\begin{aligned}
F(y_t)+E_{t+1}-E_t
&\overset{\mathrm{(A)}}{\le}
\frac{\eta D^2}{2}+L\|u_t\|-\frac{a-\beta G^2}{2}\|u_t\|^2\\
&=\frac{\eta D^2}{2}+L\|u_t\|-\frac a4\|u_t\|^2
\\&\overset{\mathrm{(B)}}{\le}\frac{\eta D^2}{2}+\frac{L^2}{a}.
\end{aligned}
\end{equation}
\nopagebreak[3]
Here (A) substitutes \eqref{eq:ch4-cb-constraint-work} into
\eqref{eq:ch4-cb-primal-comparison}, using
\eqref{eq:ch4-cb-F-definition}.
The equality uses $a=2\beta G^2$ from
\eqref{eq:ch4-cb-proof-notation}; (B) is
\eqref{eq:ch4-cb-square-bound} at $r=\|u_t\|$.

\par\medskip\noindent\textbf{Summation.}
When $T=1$, the algorithm performs no updates; sums over
$t=1,\ldots,T-1$ are empty and equal zero.
By \eqref{eq:ch4-cb-F-definition} and \eqref{eq:ch4-cb-init},
\begin{equation}\label{eq:ch4-cb-summed-comparison}
\begin{aligned}
\sum_{t=1}^{T-1}F(y_t)+E_T
&=\sum_{t=1}^{T-1}\bigl(F(y_t)+E_{t+1}-E_t\bigr)+E_1\\
&\overset{\mathrm{(A)}}{\le}
(T-1)\left(\frac{\eta D^2}{2}+\frac{L^2}{a}\right)
+\frac a2\|y_1-x^*\|^2\\
&\overset{\mathrm{(B)}}{\le}
\frac{aD^2}{2}+(T-1)\left(\frac{\eta D^2}{2}+\frac{L^2}{a}\right).
\end{aligned}
\end{equation}
\nopagebreak[3]
Here (A) sums \eqref{eq:ch4-cb-one-step}; (B) uses
$y_1=x_1\in\X$ from \eqref{eq:ch4-cb-init} and
the diameter bound in \eqref{eq:ch4-cb-domain}.

By \eqref{eq:ch4-cb-domain} and \eqref{eq:ch4-cb-clipped-zero-violation-set},
$f(x^*)=f^*$ and $c(x^*)=0$. Thus
\begin{equation}\label{eq:ch4-cb-terminal-F-bound}
\begin{aligned}
F(y_T)
&=f(y_T)-f(x^*)+\frac\beta2\|c(y_T)-c(x^*)\|^2\\
&\overset{\mathrm{(A)}}{\le}
L\|y_T-x^*\|+\frac{\beta G^2}{2}\|y_T-x^*\|^2\\
&=\frac a2\|y_T-x^*\|^2
+L\|y_T-x^*\|-\frac a4\|y_T-x^*\|^2\\
&\overset{\mathrm{(B)}}{\le}\frac a2\|y_T-x^*\|^2+\frac{L^2}{a}.
\end{aligned}
\end{equation}
\nopagebreak[3]
The equalities use the definition of $F$ and $\beta G^2=a/2$.
Here (A) uses the Lipschitz bounds
\eqref{eq:ch4-cb-scalar-lipschitz}--\eqref{eq:ch4-cb-lipschitz};
(B) is \eqref{eq:ch4-cb-square-bound} at $r=\|y_T-x^*\|$.

Consequently,
\begin{equation}\label{eq:ch4-cb-full-sum}
\begin{aligned}
&\sum_{t=1}^T F(y_t)+\frac{\|q_T\|^2}{2\eta}\\
&\quad=\sum_{t=1}^{T-1}F(y_t)+E_T
+F(y_T)-\frac a2\|y_T-x^*\|^2-\frac\beta2\|w_T\|^2\\
&\quad\overset{\mathrm{(A)}}{\le}
\sum_{t=1}^{T-1}F(y_t)+E_T+\frac{L^2}{a}\\
&\quad\overset{\mathrm{(B)}}{\le}
\frac{aD^2}{2}+\frac{(T-1)\eta D^2}{2}+\frac{TL^2}{a}\\
&\quad=LD\sqrt T+\frac{(T-1)LD}{\sqrt T}+\frac{LD\sqrt T}{2}
\overset{\mathrm{(C)}}{\le}\frac52LD\sqrt T.
\end{aligned}
\end{equation}
\nopagebreak[3]
The equalities use \eqref{eq:ch4-cb-F-definition} and
\eqref{eq:ch4-cb-tuning}, respectively.
Here (A) uses \eqref{eq:ch4-cb-terminal-F-bound} and drops
$-\beta\|w_T\|^2/2$; (B) uses \eqref{eq:ch4-cb-summed-comparison};
(C) uses $T-1\le T$.

\par\medskip\noindent\textbf{The returned point.}
By \eqref{eq:ch4-cb-init} and \eqref{eq:ch4-cb-multiplier-updates},
\begin{equation}\label{eq:ch4-cb-average-identity}
\begin{aligned}
\bar y_T&\coloneqq\frac1T\sum_{t=1}^T y_t,\qquad
r_T\coloneqq\|\bar x_T-\bar y_T\|,\\
q_T&=\eta\sum_{t=1}^{T-1}(x_{t+1}-y_{t+1})\\
&=\eta\sum_{t=1}^T(x_t-y_t)\\
&=\eta T(\bar x_T-\bar y_T),\\
\frac{\|q_T\|^2}{2\eta T}&=\frac{\eta T}{2}r_T^2=\frac a2r_T^2.
\end{aligned}
\end{equation}
The last equality uses $a=\eta T$ from
\eqref{eq:ch4-cb-proof-notation}.

Convexity gives $\bar x_T\in\X$ and $\bar y_T\in\Y$. Hence
\begin{equation}\label{eq:ch4-cb-average-bound}
\begin{aligned}
F(\bar y_T)+\frac a2r_T^2
&=F(\bar y_T)+\frac{\|q_T\|^2}{2\eta T}\\
&\overset{\mathrm{(A)}}{\le}
\frac1T\left(\sum_{t=1}^T F(y_t)+\frac{\|q_T\|^2}{2\eta}\right)
\overset{\mathrm{(B)}}{\le}\frac{5LD}{2\sqrt T}.
\end{aligned}
\end{equation}
\nopagebreak[3]
The equality is \eqref{eq:ch4-cb-average-identity};
(A) applies Jensen's inequality to the convex function $F$ defined in
\eqref{eq:ch4-cb-F-definition}; (B) divides
\eqref{eq:ch4-cb-full-sum} by $T$.

By \eqref{eq:ch4-cb-output} and \eqref{eq:ch4-cb-average-identity},
\begin{equation}\label{eq:ch4-cb-output-comparisons}
\begin{aligned}
\Delta_T&=f(\bar x_T)-f^*
\overset{\mathrm{(A)}}{\le}
f(\bar y_T)-f^*+Lr_T,\\
V_T&=\|c(\bar x_T)\|
\overset{\mathrm{(B)}}{\le}\|c(\bar y_T)\|+Gr_T.
\end{aligned}
\end{equation}
\nopagebreak[3]
Here (A) and (B) use \eqref{eq:ch4-cb-scalar-lipschitz} and
\eqref{eq:ch4-cb-v-lipschitz}, respectively.

Thus
\begin{equation}\label{eq:ch4-cb-joint-transfer}
\begin{aligned}
\Delta_T+\frac\beta4V_T^2
&\overset{\mathrm{(A)}}{\le}
f(\bar y_T)-f^*+Lr_T
+\frac\beta4\left(\|c(\bar y_T)\|+Gr_T\right)^2\\
&\overset{\mathrm{(B)}}{\le}
f(\bar y_T)-f^*+Lr_T
+\frac\beta2\|c(\bar y_T)\|^2+\frac{\beta G^2}{2}r_T^2\\
&=F(\bar y_T)+Lr_T+\frac a4r_T^2\\
&=F(\bar y_T)+\frac a2r_T^2+\left(Lr_T-\frac a4r_T^2\right)\\
&\overset{\mathrm{(C)}}{\le}
F(\bar y_T)+\frac a2r_T^2+\frac{L^2}{a}\\
&\overset{\mathrm{(D)}}{\le}\frac{5LD}{2\sqrt T}+\frac{L^2}{a}
\\&=\frac{5LD}{2\sqrt T}+\frac{LD}{2\sqrt T}
=\frac{3LD}{\sqrt T}.
\end{aligned}
\end{equation}
\nopagebreak[3]
Here (A) uses \eqref{eq:ch4-cb-output-comparisons}, (B) bounds the square
of a sum by twice the sum of the squares, (C) is \eqref{eq:ch4-cb-square-bound}
at $r=r_T$, and (D) is \eqref{eq:ch4-cb-average-bound}.
The first equality uses \eqref{eq:ch4-cb-F-definition} and
$a=2\beta G^2$ from \eqref{eq:ch4-cb-proof-notation};
the final substitution uses \eqref{eq:ch4-cb-tuning}.

Hence
\[
\frac{\Delta_T}{LD}
+\frac{\sqrt T}{4}\left(\frac{V_T}{GD}\right)^2
=\frac{\Delta_T+\beta V_T^2/4}{LD}
\overset{\mathrm{(A)}}{\le}\frac3{\sqrt T}.
\]
\nopagebreak[3]
The equality uses \eqref{eq:ch4-cb-tuning}; (A) divides
\eqref{eq:ch4-cb-joint-transfer} by $LD$.
This proves \eqref{eq:ch4-cb-theorem}; its objective upper bound
follows by dropping the nonnegative squared term.

\par\medskip\noindent\textbf{Constraint violation.}
Under Assumption~\ref{assum:Lagrange multiplier},
\begin{equation}\label{eq:ch4-cb-dual-lower-bound}
\Delta_T
\overset{\mathrm{(A)}}{\ge}-\mu^\top h(\bar x_T)
\overset{\mathrm{(B)}}{\ge}-\mu^\top c(\bar x_T)
\overset{\mathrm{(C)}}{\ge}-\|\mu\|V_T.
\end{equation}
\nopagebreak[3]
Here (A) is \eqref{eq:ch4-cb-multiplier} at $\bar x_T\in\X$,
(B) uses $\mu\ge0$ and \eqref{eq:ch4-cb-clipped-constraints};
(C) is Cauchy--Schwarz with
\eqref{eq:ch4-cb-output}.

Combining the joint objective and constraint violation bound
\eqref{eq:ch4-cb-joint-transfer} with the objective lower bound
\eqref{eq:ch4-cb-dual-lower-bound} gives
\begin{equation}\label{eq:ch4-cb-violation-quadratic}
\begin{aligned}
&V_T\left(\frac{\sqrt T\,V_T}{4GD}-\frac{G\|\mu\|}{L}\right)\\
&\quad=\frac GL\left(\frac\beta4V_T^2-\|\mu\|V_T\right)\\
&\quad\overset{\mathrm{(A)}}{\le}\frac GL\left(\frac\beta4V_T^2+\Delta_T\right)
\overset{\mathrm{(B)}}{\le}\frac{3GD}{\sqrt T}.
\end{aligned}
\end{equation}
\nopagebreak[3]
The equality uses \eqref{eq:ch4-cb-tuning}; (A) uses
\eqref{eq:ch4-cb-dual-lower-bound}, and (B) uses
\eqref{eq:ch4-cb-joint-transfer}.
\par\Needspace{5\baselineskip}
Suppose, for contradiction, that
\[
V_T\ge\frac{4GD}{\sqrt T}\left(1+\frac{G\|\mu\|}{L}\right).
\]
Then
\[
\frac{\sqrt T\,V_T}{4GD}-\frac{G\|\mu\|}{L}
\ge\left(1+\frac{G\|\mu\|}{L}\right)-\frac{G\|\mu\|}{L}
=1.
\]
Since $V_T\ge0$, \eqref{eq:ch4-cb-violation-quadratic} therefore gives
\[
V_T
\le V_T\left(\frac{\sqrt T\,V_T}{4GD}-\frac{G\|\mu\|}{L}\right)
\le\frac{3GD}{\sqrt T}.
\]
The upper bound $V_T\le3GD/\sqrt T$ contradicts the supposition
$V_T\ge(4GD/\sqrt T)(1+G\|\mu\|/L)\ge4GD/\sqrt T$.
This proves the violation bound in \eqref{eq:ch4-cb-violation-theorem}.
The lower bound $\Delta_T\ge-\|\mu\|V_T$ follows from
\eqref{eq:ch4-cb-dual-lower-bound}, and the upper bound
$\Delta_T\le3LD/\sqrt T$ follows from \eqref{eq:ch4-cb-theorem}.
\end{proof}

\begin{Thm:Corollary}[Oracle complexity]\label{cor:ch4-cb-iteration-complexity}
Under Assumptions~\ref{ass:ch2-oracles} and~\ref{assum:Lagrange multiplier},
Algorithm~\ref{alg:functional} with \eqref{eq:ch4-cb-tuning} satisfies
$\Delta_T\le\epsilon$ and $V_T\le\epsilon$ for any $\epsilon>0$
whenever its integer query budget obeys
\begin{equation}\label{eq:ch4-cb-sufficient-iteration-count}
T\ge\max\left\{1,
\left\lceil\frac{9L^2D^2}{\epsilon^2}\right\rceil,
\left\lceil\frac{16G^2D^2}{\epsilon^2}
\left(1+\frac{G\|\mu\|}{L}\right)^2\right\rceil\right\}.
\end{equation}
\end{Thm:Corollary}
\begin{proof}
The budget in \eqref{eq:ch4-cb-sufficient-iteration-count} gives
$\sqrt T\ge3LD/\epsilon$ and
$\sqrt T\ge(4GD/\epsilon)(1+G\|\mu\|/L)$.
Substituting these inequalities into
Theorem~\ref{thm:ch4-cb-guarantee} proves both claims.
\end{proof}

\section{Conclusion}

Using linear optimization over $\X$ and projection onto the auxiliary
set $\Y$, the method returns a point with objective gap at most
$3LD/\sqrt T$ by~\eqref{eq:ch4-cb-theorem}. This bound is computable from
the known constants, and all update coefficients depend only on the
query budget $T$ and the bounds $D,L,G$.

Under Assumption~\ref{assum:Lagrange multiplier}, the constraint violation
norm also has an $O(T^{-1/2})$ upper bound. Evaluating this theoretical
bound in \eqref{eq:ch4-cb-violation-theorem} requires a known constant
$M\ge\|\mu\|$ for a multiplier $\mu$ satisfying
\eqref{eq:ch4-cb-multiplier}. The returned point's violation can also be
checked by evaluating $h(\bar x_T)$. Unless its value is already
available, this check requires one additional constraint-vector
evaluation beyond the stated optimization budget.
